\section*{Bab \ref{ch:g11:antibiotic-resistance} --- Bakteri dan Ketahanan terhadap Antibiotik}

\begin{solution}{exo:g11:antibiotic-resistance:1}
\omterm{def:g11:mutations:mutation}{Mutasi} yang muncul selama pembelahan klonnya, dan \omterm{prop:g11:antibiotic-resistance:variation}{pemindahan gen
mendatar} --- yaitu \omterm{def:g10:universal-dna:information}{DNA}, biasanya sebuah plasmid, yang diterima dari
bakteri lain.
\end{solution}

\begin{solution}{exo:g11:antibiotic-resistance:2}
Molekul yang membunuh bakteri atau menghentikan pertumbuhannya dengan
bekerja pada struktur yang dimiliki bakteri dan tidak dimiliki \omterm{def:g10:cells-common-unit:cell}{sel}
manusia (dinding, \omterm{def:g10:cells-common-unit:organelle}{ribosom} bakteri, \omterm{def:g11:enzymes-and-phenotype:enzyme}{enzim} replikasi bakteri). Pilek
disebabkan virus, dan virus tidak punya satu pun sasaran itu.
\end{solution}

\begin{solution}{exo:g11:antibiotic-resistance:3}
\omterm{def:g11:enzymes-and-phenotype:enzyme}{Enzim} yang memusnahkan \omterm{def:g11:antibiotic-resistance:antibiotic}{antibiotiknya}; sasaran yang bermutasi yang tidak
lagi diikat obatnya; pompa yang mengusir obatnya (atau dinding yang
menolaknya).
\end{solution}

\begin{solution}{exo:g11:antibiotic-resistance:4}
Peka (20 di atas 18), tetapi hanya sedikit. D, yang selisihnya paling
lebar di atas ambangnya, lalu A.
\end{solution}

\begin{solution}{exo:g11:antibiotic-resistance:5}
Tidak. Pada cawan replikanya, koloni yang \omterm{prop:g11:antibiotic-resistance:mechanisms}{tahan} muncul pada kedudukan
yang sama pada tiap cawan \omterm{def:g11:antibiotic-resistance:antibiotic}{antibiotiknya}, sehingga \omterm{def:g10:cells-common-unit:cell}{sel} yang \omterm{prop:g11:antibiotic-resistance:mechanisms}{tahan} sudah
ada dalam koloni aslinya, yang tidak pernah bertemu obatnya.
\omterm{def:g11:antibiotic-resistance:antibiotic}{Antibiotiknya} menyingkapkan dan memilah; ia tidak mengimbaskan.
\end{solution}

\begin{solution}{exo:g11:antibiotic-resistance:6}
30 pembelahan: $2^{30} \approx 10^9$ \omterm{def:g10:cells-common-unit:cell}{sel}. Sekitar $10^9$ pembelahan
menghasilkannya: sekitar 10 \omterm{def:g10:cells-common-unit:cell}{sel} yang \omterm{prop:g11:antibiotic-resistance:mechanisms}{tahan}, rata-rata --- lebih banyak
jika \omterm{def:g11:mutations:mutation}{mutasinya} datang lebih awal.
\end{solution}

\begin{solution}{exo:g11:antibiotic-resistance:7}
Sekitar hari ke-3, ketika \omterm{def:g10:cells-common-unit:cell}{sel} yang peka telah turun menjadi sekitar
0.2\% dan yang \omterm{prop:g11:antibiotic-resistance:mechanisms}{tahan} telah naik menjadi sekitar 1\%. Pada hari ke-6
jumlah bergaris putus-putusnya kembali pada 100\%, yang hampir
seluruhnya tersusun dari \omterm{def:g10:cells-common-unit:cell}{sel} yang \omterm{prop:g11:antibiotic-resistance:mechanisms}{tahan} (\omterm{def:g10:cells-common-unit:cell}{sel} peka yang selamat pada
hari ke-3 juga tumbuh kembali, tetapi dari alas yang lebih kecil).
\end{solution}

\begin{solution}{exo:g11:antibiotic-resistance:8}
\omterm{def:g10:universal-dna:gene}{Gen} \omterm{def:g11:cell-cycle-mitosis:chromosome}{kromosom} hanya berpindah kepada keturunan klonnya; sebuah plasmid
juga berpindah secara mendatar, kepada \omterm{def:g10:cells-common-unit:cell}{sel} yang tak berkerabat dan
\omterm{def:g10:biodiversity-scales:species}{spesies} lain, dan satu pemindahan memberi sebuah \omterm{def:g10:cells-common-unit:cell}{sel} ketahanan seketika
tanpa menunggu sebuah \omterm{def:g11:mutations:mutation}{mutasi}.
\end{solution}

\begin{solution}{exo:g11:antibiotic-resistance:9}
Pada hari ke-3 \omterm{def:g10:cells-common-unit:cell}{sel} yang peka telah turun ke sepersekian persen
sedangkan \omterm{def:g10:cells-common-unit:cell}{sel} yang \omterm{prop:g11:antibiotic-resistance:mechanisms}{tahan} sedang naik; tanpa obatnya, seluruh yang
selamat tumbuh kembali, dan populasi yang kembali jauh lebih kaya \omterm{def:g10:cells-common-unit:cell}{sel}
yang \omterm{prop:g11:antibiotic-resistance:mechanisms}{tahan} daripada populasi aslinya. Kambuhnya tidak lagi menanggapi
obat yang sama.
\end{solution}

\begin{solution}{exo:g11:antibiotic-resistance:10}
Rumah sakit memakai \omterm{def:g11:antibiotic-resistance:antibiotic}{antibiotik} secara giat pada banyak pasien yang
berdekatan: seleksinya tetap dan galur yang terpilih berpindah dari
pasien ke pasien lewat tangan dan permukaan.
\end{solution}

\begin{solution}{exo:g11:antibiotic-resistance:11}
Streptomisin bekerja dengan mengikat sebuah \omterm{def:g11:gene-expression:protein}{protein} \omterm{def:g10:cells-common-unit:organelle}{ribosom}; \omterm{def:g11:gene-expression:protein}{protein}
yang berubah tidak lagi mengikatnya dan translasinya berlanjut. Tetapi
perubahannya juga sedikit melemahkan kerja normal \omterm{def:g10:cells-common-unit:organelle}{ribosomnya}, sehingga
penyusunan \omterm{def:g11:gene-expression:protein}{proteinnya}, dan pertumbuhannya, sedikit lebih lambat ---
yaitu harga ketahanannya.
\end{solution}

\begin{solution}{exo:g11:antibiotic-resistance:12}
$10^{-8} \times 10^{-8} = 10^{-16}$ tiap pembelahan: dengan $10^{12}$
bakteri, tidak mungkin ada \omterm{def:g10:cells-common-unit:cell}{sel} yang membawa keduanya. Pada
tuberkulosis, pengobatannya berlangsung berbulan-bulan dan bakterinya
bermiliar-miliar, sehingga ketahanan terhadap satu obat pasti terpilih;
dua obat sekaligus tidak meninggalkan yang selamat untuk dipilah.
\end{solution}

\begin{solution}{exo:g11:antibiotic-resistance:13}
Dosis rendah memilih galur yang \omterm{prop:g11:antibiotic-resistance:mechanisms}{tahan} dalam usus ternaknya tanpa
membunuh populasinya; pekerjanya memperolehnya lewat sentuhan;
plasmidnya membawa \omterm{def:g10:universal-dna:gene}{gennya} ke dalam kuman penyakit manusia dan ke
lingkungan lewat pupuk kandang. Ketahanan terhadap obat itu, dan
terhadap obat manusia yang berkerabat, naik dalam kedokteran.
\end{solution}

\begin{solution}{exo:g11:antibiotic-resistance:14}
Ketahanannya sering menuntut harga dari bakterinya (\omterm{def:g10:cells-common-unit:organelle}{ribosom} yang lebih
lambat, \omterm{def:g11:enzymes-and-phenotype:enzyme}{enzim} yang harus dibuat): tanpa obatnya, \omterm{def:g10:cells-common-unit:cell}{sel} yang peka tumbuh
sedikit lebih cepat dan mengalahkan yang \omterm{prop:g11:antibiotic-resistance:mechanisms}{tahan} sepanjang banyak
keturunan. Bagiannya turun --- tetapi \omterm{def:g10:universal-dna:gene}{gen} yang \omterm{prop:g11:antibiotic-resistance:mechanisms}{tahan} tetap ada dalam
populasinya pada frekuensi yang rendah, siap dipilih lagi.
\end{solution}

\begin{solution}{exo:g11:antibiotic-resistance:15}
Bakteri tidak menyesuaikan diri sendiri-sendiri. Dalam populasi besar
mana pun, beberapa \omterm{def:g10:cells-common-unit:cell}{selnya} membawa, lewat \omterm{def:g11:mutations:mutation}{mutasi} acak atau lewat plasmid
yang diterima dari \omterm{def:g10:cells-common-unit:cell}{sel} lain, sebuah \omterm{def:g10:universal-dna:gene}{alel} yang membuatnya \omterm{prop:g11:antibiotic-resistance:mechanisms}{tahan}.
\omterm{def:g11:antibiotic-resistance:antibiotic}{Antibiotiknya} membunuh yang lain; \omterm{def:g10:cells-common-unit:cell}{sel} yang \omterm{prop:g11:antibiotic-resistance:mechanisms}{tahan} berlipat ganda dan
mewariskan \omterm{def:g10:universal-dna:gene}{alelnya} kepada keturunannya. Populasinya menjadi \omterm{prop:g11:antibiotic-resistance:mechanisms}{tahan}
karena ia telah terpilah, bukan karena anggotanya belajar sesuatu.
\end{solution}

\begin{solution}{pb:g11:antibiotic-resistance:1}
\textbf{1.} Sekitar $10^{10}$ pembelahan (satu tiap \omterm{def:g10:cells-common-unit:cell}{sel} yang dihasilkan).

\textbf{2.} $10^{10} \times 10^{-8} = 100$ \omterm{def:g10:cells-common-unit:cell}{sel} yang \omterm{prop:g11:antibiotic-resistance:mechanisms}{tahan}, rata-rata.

\textbf{3.} Seratus \omterm{def:g10:cells-common-unit:cell}{sel} dalam sepuluh miliar tak terlihat pada sebuah
cawan: hamparannya berperilaku peka karena 99.999999\% darinya peka.

\textbf{4.} Amoksisilin: peka (24 di atas 17). Eritromisin: \omterm{prop:g11:antibiotic-resistance:mechanisms}{tahan} (12
di bawah 20). Penisilin: \omterm{prop:g11:antibiotic-resistance:mechanisms}{tahan} sepenuhnya, tanpa hambatan.

\textbf{5.} Penisilin tidak berguna melawan galur ini. Pilihan di
antara obat yang peka juga menimbang efek samping dan cakupannya; obat
tersempit yang berhasil lebih disukai agar bakteri pasiennya yang lain
tersayangi.

\textbf{6.} Empat kurun 6 jam tiap hari: $10^{10} \times 0.1^4 =
10^6$ setelah 1 hari; $10^{10} \times 0.1^{12} = 0.01$ setelah 3 hari
--- pada dasarnya tidak ada.

\textbf{7.} $100 \times 2^{48} \approx \num{2.8e16}$: mustahil, lebih
banyak daripada yang dapat dimuat lukanya. Pertumbuhannya dibatasi
ruang, zat gizi dan sistem kekebalannya.

\textbf{8.} Dari 100 \omterm{def:g10:cells-common-unit:cell}{sel} menjadi $10^{10}$ menuntut $\log_2 10^8 \approx 27$
pelipatduaan, sekitar 13 jam; dalam praktiknya ruangnya bebas selama
hari pertamanya, sehingga pada hari ke-1 sampai ke-2 klon yang \omterm{prop:g11:antibiotic-resistance:mechanisms}{tahan}
telah menggantikan populasi yang peka.

\textbf{9.} Yang peka: turun dengan faktor 10 tiap 6 jam sampai tidak
ada pada hari ke-3. Yang \omterm{prop:g11:antibiotic-resistance:mechanisms}{tahan}: naik dari 100 untuk mengisi ruangnya
dalam sehari atau dua hari, lalu mendatar pada $10^{10}$. Seluruhnya:
sebuah cekungan selama hari ke-1, lalu kembali ke $10^{10}$, yang kini
seluruhnya \omterm{prop:g11:antibiotic-resistance:mechanisms}{tahan}.

\textbf{10.} Jika sistem kekebalannya membersihkan \omterm{def:g10:cells-common-unit:cell}{sel} lebih cepat
daripada seratus \omterm{def:g10:cells-common-unit:cell}{sel} \omterm{prop:g11:antibiotic-resistance:mechanisms}{tahan} itu dapat berlipat ganda, infeksinya
berakhir walaupun \omterm{def:g10:cells-common-unit:cell}{sel} itu ada; jika klonnya tumbuh melampaui
pembersihannya, infeksinya menjadi \omterm{prop:g11:antibiotic-resistance:mechanisms}{tahan}. \omterm{def:g11:antibiotic-resistance:antibiotic}{Antibiotik} bekerja bersama
sistem kekebalannya, bukan menggantikannya.

\textbf{11.} Yang peka setelah 2 hari: $10^{10} \times 0.1^8 = 100$.
Yang \omterm{prop:g11:antibiotic-resistance:mechanisms}{tahan}: sekitar $10^{10}$ jika klonnya telah mengisi ruangnya, atau
sedikitnya beberapa ribu: yang selamat hampir seluruhnya \omterm{prop:g11:antibiotic-resistance:mechanisms}{tahan}.

\textbf{12.} Tumbuh dengan laju yang sama menjaga perbandingannya: pada
dasarnya 100\% \omterm{prop:g11:antibiotic-resistance:mechanisms}{tahan}.

\textbf{13.} Tanpa pengaruh: klon yang \omterm{prop:g11:antibiotic-resistance:mechanisms}{tahan} mengabaikan amoksisilin.

\textbf{14.} Cawan pertamanya memuat populasi aslinya, yang
99.999999\% peka; cawan kambuhnya memuat klon yang terpilih, yang
seluruhnya \omterm{prop:g11:antibiotic-resistance:mechanisms}{tahan}, sehingga hamparannya tumbuh sampai ke cakramnya.

\textbf{15.} \omterm{def:g10:universal-dna:gene}{Gennya} dapat berpindah kepada \omterm{def:g10:biodiversity-scales:species}{spesies} lain yang dibawa
pasien atau petugas lain, sehingga kuman penyakit lain di bangsalnya
dapat menjadi \omterm{prop:g11:antibiotic-resistance:mechanisms}{tahan} tanpa menunggu \omterm{def:g11:mutations:mutation}{mutasinya} sendiri.

\textbf{16.} Di bangsalnya, 400 rangkaian tiap tahun dalam populasi
kecil yang tertutup memilah bakterinya terus-menerus, dan galur yang
terpilih berpindah antarpasien; di kotanya seleksinya terencerkan di
antara banyak orang dan \omterm{def:g10:biodiversity-scales:species}{spesies}.

\textbf{17.} Antibiogram: pilihlah obat yang tidak dapat ditahan
galurnya, sehingga tidak ada klon yang terpilih. Rangkaian yang penuh:
tidak ada rangkaian terputus yang meninggalkan yang \omterm{prop:g11:antibiotic-resistance:mechanisms}{tahan} untuk tumbuh
kembali. Kebersihan: tidak ada penularan klon yang terpilih
antarpasien. Pengasingan: tidak ada penyebaran mendatar atau tegak dari
\omterm{def:g11:genes-and-disease:genetic}{pembawanya}.

\textbf{18.} Seleksi yang lebih sedikit membiarkan galur yang peka,
yang tumbuh sedikit lebih cepat, merebut kembali tempatnya, dan
penularan galur yang \omterm{prop:g11:antibiotic-resistance:mechanisms}{tahan} terpotong. Ia tidak kembali ke 5\% karena
galur yang \omterm{prop:g11:antibiotic-resistance:mechanisms}{tahan} dan plasmidnya bertahan pada frekuensi yang rendah dan
terpilih lagi oleh tiap rangkaian yang tersisa.

\textbf{19.} Populasi bakteri besar mana pun memuat \omterm{def:g10:cells-common-unit:cell}{sel} yang \omterm{prop:g11:antibiotic-resistance:mechanisms}{tahan}
terhadap obat barunya lewat \omterm{def:g11:mutations:mutation}{mutasi} kebetulan; memakainya memilih \omterm{def:g10:cells-common-unit:cell}{sel}
itu, dan dalam beberapa tahun galur yang \omterm{prop:g11:antibiotic-resistance:mechanisms}{tahan} menjadi lazim, seperti
pada setiap \omterm{def:g11:antibiotic-resistance:antibiotic}{antibiotik} yang diperkenalkan sejauh ini.

\textbf{20.} Sekitar seratus \omterm{def:g10:cells-common-unit:cell}{sel} yang \omterm{prop:g11:antibiotic-resistance:mechanisms}{tahan} di antara sepuluh miliar
sebelum obat apa pun; pada hari kedua pengobatannya \omterm{def:g10:cells-common-unit:cell}{sel} itu mewarisi
lukanya; menghentikan rangkaiannya lebih awal mengubah infeksi yang
dapat disembuhkan menjadi infeksi yang \omterm{prop:g11:antibiotic-resistance:mechanisms}{tahan} dan menyebar.
\end{solution}
